\paragraph{Task vectors and per-task geometry.} Fix integers $T \geq 1$,
$d \geq 2$, $1 \leq r \leq d$, and a radius $B > 0$. Each of the $T$
tasks is specified by a pair $(\tau_t, H_t) \in \R^d \times \R^{d\times d}$
with
\begin{equation}\label{eq:admissibility}
  H_t \succeq 0,\quad \rank(H_t) \leq r,\quad \tau_t \in V_t :=
  \range(H_t),\quad \tau_t^\top H_t \tau_t \leq B^2.
\end{equation}
Per-task loss is the quadratic $f_t(w) := \tfrac12 (w - \tau_t)^\top
H_t (w - \tau_t)$. We write $\boldsymbol\tau := (\tau_t)_{t=1}^T$,
$\boldsymbol H := (H_t)_{t=1}^T$, and call tuples satisfying
\eqref{eq:admissibility} \emph{admissible}. The specialization
$H_t = P_{V_t}$ recovers the weight-disentangled LoRA setting
of~\citet{ortizjimenez2023disentanglement}; general $H_t \succeq 0$
subsumes Fisher-metric surrogates for cross-entropy.

\paragraph{Worst-task distortion.} A merged model $w^\star \in \R^d$
incurs per-task distortion $D_t(w^\star) := (w^\star - \tau_t)^\top H_t
(w^\star - \tau_t)$; the figure of merit is the \emph{worst-task}
distortion $\Delta(w^\star; \boldsymbol\tau, \boldsymbol H) :=
\max_{t\in[T]} D_t(w^\star)$: a merged model that collapses on one
task is not deployable, and classical multi-description
coding~\citep{elgamalcover1982md} treats average distortion, which has
a strictly weaker rate-distortion function (Lemma~\ref{lem:id}).

\paragraph{Rate-$R$ merging codes.} A rate-$R$ merging code is a pair
$(E, D)$ with encoder $E : (\R^d)^T \to \{1,\ldots,2^R\}$ and decoder
$D : \{1,\ldots,2^R\} \to \R^d$, both of which may depend on
$\boldsymbol H$ as shared side information: in deployment
$\boldsymbol H$ is determined by the LoRA factors both ends already
hold, so only the $R$ bits describing $\boldsymbol\tau$ are charged to
the rate. Our lower bound holds even when $\boldsymbol H$ is revealed
for free to both ends. The \emph{rank-$r$ rate-distortion function} is
\begin{equation}\label{eq:rd-def}
  \mathcal{D}^\star(T, d, r, B, R)
  \;:=\; \inf_{(E, D)}\;\sup_{(\boldsymbol\tau, \boldsymbol H)
    \text{ admissible}}\;
  \E\bigl[\Delta(w^\star;\boldsymbol\tau, \boldsymbol H)\bigr],
  \qquad w^\star := D(E(\boldsymbol\tau)),
\end{equation}
with the expectation over any shared randomness. In the floor-zero
Stiefel regime the paper determines $\mathcal{D}^\star$ up to a
universal constant: a lower bound for every code
(Theorem~\ref{thm:general}) and an explicit code matching it
(Theorem~\ref{thm:achv}).

\paragraph{Key derived quantities.} We repeatedly use
\begin{equation}\label{eq:derived}
  \bar H := \tfrac1T\sum_{t=1}^T H_t,\qquad
  \tauH := \bar H^+\cdot \tfrac1T\sum_{t=1}^T H_t\tau_t,\qquad
  \deff := \rank\!\Bigl(\sum_{t=1}^T P_{V_t}\Bigr),
\end{equation}
where $(\cdot)^+$ is the Moore--Penrose pseudoinverse. The
\emph{$H_t$-weighted centroid} $\tauH$ is the unique element of
$\range(\bar H) = V_1 + \cdots + V_T$ with $\bar H \tauH =
T^{-1}\sum_t H_t \tau_t$, reducing to the arithmetic mean when
$H_t = I_d$; the \emph{effective dimension} $\deff \leq \min(Tr, d)$
counts the independent directions the decoder must reconstruct.

\paragraph{Hard distribution $P^\star$.} The lower bound is established
against a distribution $P^\star$ on admissible tuples. For each task
$t$, independently: (i) sample $U_t \in \mathrm{Stiefel}(d, r)$
uniformly and set $V_t := \range(U_t)$; (ii) fix a positive-definite
diagonal $D_t \succ 0$ (the task-$t$ spectrum) and set $H_t := U_t D_t
U_t^\top$; (iii) sample $u_t \sim \mathrm{Unif}(S^{r-1})$ and set
$\tau_t := U_t\, B D_t^{-1/2} u_t$. The normalization
$\tau_t^\top H_t \tau_t = B^2$ holds by construction, and the
$D_t^{-1/2}$ stretch makes the second-moment identity
$\E[H_t \tau_t \tau_t^\top H_t \mid U_t] = (B^2/r) H_t$ hold
\emph{regardless of} $D_t$, the feature that makes
Theorem~\ref{thm:general} apply to arbitrary task-dependent spectra.

\paragraph{Yao's minimax bridge.} By Yao's minimax
principle~\citep{yao1977}, $\mathcal{D}^\star$ is at least the
Bayesian rate-distortion function under the prior $P^\star$:
\begin{equation}\label{eq:yao}
  \mathcal{D}^\star(T,d,r,B,R)
  \;\geq\;
  \inf_{(E,D)}\;\E_{(\boldsymbol\tau,\boldsymbol H)\sim P^\star}
    \bigl[\Delta(w^\star;\boldsymbol\tau,\boldsymbol H)\bigr]
\end{equation}
All lower bounds below are proved against the right-hand side.
